% Part: many-valued-logic
% Chapter: syntax-and-semantics
% Section: matrices

\documentclass[../../../include/open-logic-section]{subfiles}

\begin{document}

\olfileid[fr]{mvl}{syn}{mat}

\olsection{Matrices}

Une logique plurivalente est donnée
par son langage, son ensemble~$V$ de
valeurs de vérité, un sous-ensemble
de valeurs désignées et les fonctions
de vérité de ses connecteurs. Ces
données forment une \emph{matrice}.

\begin{defn}[Matrice]
\ollabel{defn:matrix}
Une \emph{matrice} pour la logique
$\Log L$ comprend :
\begin{enumerate}
\item un ensemble de connecteurs
  constituant un langage~$\Lang L$;
\item un ensemble non vide~$V$
  de valeurs de vérité;
\item un ensemble~$V^+\subseteq V$
  de valeurs de vérité désignées;
\item pour chaque connecteur~$\star$
  d'arité~$n$ de~$\Lang L$, une
  fonction de vérité
  $\tf{\star}\colon V^n\to V$.
  Si $n=0$, $\tf{\star}$ est
  simplement un élément de~$V$.
\end{enumerate}
\end{defn}

\begin{ex}
La matrice de la logique classique
$\LogCL{}$ se compose :
\begin{enumerate}
\item du langage propositionnel usuel
  $\Lang L_0$, avec $\lfalse$,
  $\lnot$, $\land$, $\lor$ et $\lif$;
\item de l'ensemble des valeurs
  $V=\{\True,\False\}$;
\item de la seule valeur désignée
  $\True$, c'est-à-dire
  $V^+=\{\True\}$;
\item de la valeur
  $\tf{\lfalse}=\False$ et des
  fonctions de vérité des autres
  connecteurs données par les tables
  usuelles (voir la \olref{fig:tf-CL}).
\end{enumerate}
\begin{figure}
  \begin{center}
      \begin{tabular}{c|c}
        $\tf{\lnot}$ & \\
        \hline
        $\True$ & $\False$ \\
        $\False$ & $\True$
      \end{tabular}
      \quad
      \begin{tabular}{c|cc}
        $\tf{\land}$ & $\True$ & $\False$ \\
        \hline
        $\True$ & $\True$ & $\False$ \\
        $\False$ & $\False$ & $\False$
      \end{tabular}
      \quad
      \begin{tabular}{c|cc}
        $\tf{\lor}$ & $\True$ & $\False$ \\
        \hline
        $\True$ & $\True$ & $\True$ \\
        $\False$ & $\True$ & $\False$
      \end{tabular}
      \quad
      \begin{tabular}{c|cc}
        $\tf{\lif}$ & $\True$ & $\False$ \\
        \hline
        $\True$ & $\True$ & $\False$ \\
        $\False$ & $\True$ & $\True$
      \end{tabular}
    \end{center}
    \caption{Fonctions de vérité de la logique
      classique~$\LogCL$.}
    \ollabel{fig:tf-CL}
  \end{figure}
\end{ex}

\end{document}
